\section{New design guide}
As Xamarin was changed to flutter, the entire frontend had to be reworked. 
Therefor we took a look at the design guide and decided that there are things that we wanted to changed to get a better design.
This meant that a new design guide had to be worked on.
One of the things we considered were that an 80 page design guide was impractical. 
The new designguide is therefore on the github wiki alongside the rest of the GIRAF project and supposed to be easier to look up things.
As seen on \autoref{fig:designguides} the new one has fewer options, where the old one had 4 pages of content in the PDF file.
\begin{figure}[H]
    \begin{subfigure}{0.6\textwidth}
    \includegraphics[width=1\linewidth, height=5cm]{figures/table-of-content-designguide.JPG}
    \caption{New design guide table of content}
    \label{fig:new-designguide}
    \end{subfigure}
    \begin{subfigure}{0.6\textwidth}
        \includegraphics[width=1\linewidth, height=5cm]{figures/old-design-guide}
    \caption{Part of old design guide contents. Continues for another 3 pages}
    \label{fig:old-designguide}
    \end{subfigure} 
    \caption{New and old design guide.}
    \label{fig:designguides}
\end{figure} 

\subsubsection{Colors}
For the colors in GIRAF not much changed from the old designguide to the new one. 
The only thing that changed was the colors for weekday.
The new colors are less strong compared to the old colors.

\subsubsection{Typography}
For the typography the standard text color is black. The font was changed to "Quicksand font", because the customers prefers the " a " to be the same type as they learn it in school.
The size for the font is 30 pt in landscape mode and in portrait mode it should be 20 pt.
The text should also never be bold, italic or underlined to ensure consistent design.

\subsubsection{Icons}
Every icon in the old design guide has been changed for the new design guide.
The old icons had an out of date look, and didn't give a good design for the application.
Most of the icons that are used in this project are icons from "Font Awesome".
These are easy to import in flutter and by using the icons in the same package, the application is more consistent and will give the users a better experience, as they know the icons throughout the program.
Two icons that was created can be seen on \autoref{fig:change-profile-icons}. 
The new icon for changing from guardian to citizen and vice versa are icons for themselves and this will give the user an indication for which mode the user is in.
Even though this is not the only way of knowing which mode the user is in, this is an additional indicator.

\begin{figure}[htp]

    \centering
    \includegraphics[width=.1\textwidth]{figures/changeToGuardian}\hfill
    \includegraphics[width=.1\textwidth]{figures/changeToCitizen}\hfill
    \includegraphics[width=.1\textwidth]{figures/old-change-profile}
    
    \caption{The rightmost picture is the old icon for changing profile. The middle is change from guardian to citizen and the leftmost is change from citizen to guardian.}
    \label{fig:change-profile-icons}
\end{figure}


